% GENERATED FILE. Edit canonical lessons, then run npm run generate:books.
% canonical-source-hash: fnv1a64:a857628afe4088ba
% canonical-lessons: FR-C481-entorse, FR-C481-tordre, FR-C481-enfler, FR-C481-attelle, FR-C481-kine, FR-R481-revision-474-481

\chapter{Une mauvaise entorse}
\label{ch:fr-une-mauvaise-entorse}
\hlchaptermodality{481}

\begin{chapteropening}
\textbf{By the end of this chapter:} \emph{I can talk about a sprain and how it heals.}

Complete the last lesson of chapter 481: I can talk about a sprain and how it heals.
\end{chapteropening}

\section[l'entorse]{l'entorse — sprain}

\label{lesson:FR-C481-entorse}



\begin{quote}
Before the new one: say the French for to limp, then the French for bone.
\end{quote}

\subsection*{You'll want to know: l'entorse}
\textbf{l'entorse} — \textquotedblleft{}sprain\textquotedblright{}.

\begin{quote}
Il s'est fait une entorse à la cheville au tennis.
\end{quote}

\begin{quote}
\emph{He sprained his ankle playing tennis.}
\end{quote}

Feminine. \emph{Une entorse au règlement} — a breach of the rules.

\begin{cousinweb}
\emph{l'entorse} comes from the old verb \emph{entordre} \textquotedblleft{}to twist\textquotedblright{}, built from \emph{en-} and \emph{tordre}, which is inherited from Latin \emph{torquēre} \textquotedblleft{}to twist\textquotedblright{} (Lewis \& Short).
\end{cousinweb}

\begin{grammarlens}[title={What you've built}]
One more word for telling, explaining and describing at B1.
\end{grammarlens}

\subsection*{Your turn}
\begin{itemize}
\raggedright
  \item \emph{Say it:} \emph{l'entorse}
  \item \emph{Say it:} \emph{l'entorse}, once more
  \item \emph{Say it:} a sentence of your own with \emph{l'entorse}
  \item \emph{Recall:} say the French for to limp, then the French for bone, then say \emph{l'entorse} again
\end{itemize}

\begin{tcolorbox}[breakable,colback=teal!4,colframe=teal!35!black,title={Before you move on}]
What does l'entorse mean? (\textquotedblleft{}Sprain\textquotedblright{}.) Say it once more.
\end{tcolorbox}
\section[tordre]{tordre — to twist, to wring; (se tordre) to twist (a joint)}

\label{lesson:FR-C481-tordre}



\begin{quote}
Before the new one: say the French for bone, then the French for sprain.
\end{quote}

\subsection*{You'll want to know: tordre}
\textbf{tordre} — \textquotedblleft{}to twist, to wring; (se tordre) to twist (a joint)\textquotedblright{}.

\begin{quote}
Je me suis tordu le pied dans l'escalier.
\end{quote}

\begin{quote}
\emph{I twisted my foot on the stairs.}
\end{quote}

For a joint it is reflexive: \emph{se tordre la cheville}. Present \emph{je tords, il tord}; past \emph{j'ai tordu}. \emph{Se tordre de rire} — to double up laughing.

\begin{cousinweb}
\emph{tordre} is inherited from Vulgar Latin \emph{torcere}, a form of Latin \emph{torquēre} \textquotedblleft{}to twist\textquotedblright{} (Lewis \& Short).
\end{cousinweb}

\begin{grammarlens}[title={What you've built}]
One more word for telling, explaining and describing at B1.
\end{grammarlens}

\subsection*{Your turn}
\begin{itemize}
\raggedright
  \item \emph{Say it:} \emph{tordre}
  \item \emph{Say it:} \emph{tordre}, once more
  \item \emph{Say it:} a sentence of your own with \emph{tordre}
  \item \emph{Recall:} say the French for bone, then the French for sprain, then say \emph{tordre} again
\end{itemize}

\begin{tcolorbox}[breakable,colback=teal!4,colframe=teal!35!black,title={Before you move on}]
What does tordre mean? (\textquotedblleft{}To twist, to wring; (se tordre) to twist (a joint)\textquotedblright{}.) Say it once more.
\end{tcolorbox}
\section[enfler]{enfler — to swell}

\label{lesson:FR-C481-enfler}



\begin{quote}
Before the new one: say the French for sprain, then the French for to twist.
\end{quote}

\subsection*{You'll want to know: enfler}
\textbf{enfler} — \textquotedblleft{}to swell\textquotedblright{}.

\begin{quote}
Sa cheville a enflé pendant la nuit.
\end{quote}

\begin{quote}
\emph{Her ankle swelled during the night.}
\end{quote}

\emph{Enflé} — swollen. The everyday word is often \emph{gonfler}.

\begin{cousinweb}
\emph{enfler} is inherited from Latin \emph{īnflāre} \textquotedblleft{}to blow into, to puff up\textquotedblright{} (EtymDB), built on \emph{flāre} \textquotedblleft{}to blow\textquotedblright{} (Lewis \& Short).
\end{cousinweb}

\begin{grammarlens}[title={What you've built}]
One more word for telling, explaining and describing at B1.
\end{grammarlens}

\subsection*{Your turn}
\begin{itemize}
\raggedright
  \item \emph{Say it:} \emph{enfler}
  \item \emph{Say it:} \emph{enfler}, once more
  \item \emph{Say it:} a sentence of your own with \emph{enfler}
  \item \emph{Recall:} say the French for sprain, then the French for to twist, then say \emph{enfler} again
\end{itemize}

\begin{tcolorbox}[breakable,colback=teal!4,colframe=teal!35!black,title={Before you move on}]
What does enfler mean? (\textquotedblleft{}To swell\textquotedblright{}.) Say it once more.
\end{tcolorbox}
\section[l'attelle]{l'attelle — splint}

\label{lesson:FR-C481-attelle}



\begin{quote}
Before the new one: say the French for to twist, then the French for to swell.
\end{quote}

\subsection*{You'll want to know: l'attelle}
\textbf{l'attelle} — \textquotedblleft{}splint\textquotedblright{}.

\begin{quote}
Le médecin lui a mis une attelle au doigt.
\end{quote}

\begin{quote}
\emph{The doctor put a splint on his finger.}
\end{quote}

Feminine. Lighter than \emph{un plâtre}.

\begin{cousinweb}
\emph{l'attelle} is inherited from Vulgar Latin \emph{astella} (EtymDB), from Latin \emph{astula} \textquotedblleft{}a splinter, a chip of wood\textquotedblright{} (Lewis \& Short).
\end{cousinweb}

\begin{grammarlens}[title={What you've built}]
One more word for telling, explaining and describing at B1.
\end{grammarlens}

\subsection*{Your turn}
\begin{itemize}
\raggedright
  \item \emph{Say it:} \emph{l'attelle}
  \item \emph{Say it:} \emph{l'attelle}, once more
  \item \emph{Say it:} a sentence of your own with \emph{l'attelle}
  \item \emph{Recall:} say the French for to twist, then the French for to swell, then say \emph{l'attelle} again
\end{itemize}

\begin{tcolorbox}[breakable,colback=teal!4,colframe=teal!35!black,title={Before you move on}]
What does l'attelle mean? (\textquotedblleft{}Splint\textquotedblright{}.) Say it once more.
\end{tcolorbox}
\section[le kiné]{le kiné — physiotherapist (informal)}

\label{lesson:FR-C481-kine}



\begin{quote}
Before the new one: say the French for to swell, then the French for splint.
\end{quote}

\subsection*{You'll want to know: le kiné}
\textbf{le kiné} — \textquotedblleft{}physiotherapist (informal)\textquotedblright{}.

\begin{quote}
Je vais chez le kiné deux fois par semaine.
\end{quote}

\begin{quote}
\emph{I go to the physio twice a week.}
\end{quote}

Short for \emph{le kinésithérapeute}. A woman is \emph{la kiné}. In Belgium also \emph{le kiné}; in Québec \emph{le physiothérapeute}.

\begin{cousinweb}
\emph{le kiné} is short for \emph{kinésithérapeute}, a learned word built from Greek \emph{kínēsis} \textquotedblleft{}movement\textquotedblright{} and \emph{therapeut\'{\={e}}s} \textquotedblleft{}one who treats\textquotedblright{}.
\end{cousinweb}

\begin{grammarlens}[title={What you've built}]
One more word for telling, explaining and describing at B1.
\end{grammarlens}

\subsection*{Your turn}
\begin{itemize}
\raggedright
  \item \emph{Say it:} \emph{le kiné}
  \item \emph{Say it:} \emph{le kiné}, once more
  \item \emph{Say it:} a sentence of your own with \emph{le kiné}
  \item \emph{Recall:} say the French for to swell, then the French for splint, then say \emph{le kiné} again
\end{itemize}

\begin{tcolorbox}[breakable,colback=teal!4,colframe=teal!35!black,title={Before you move on}]
What does le kiné mean? (\textquotedblleft{}Physiotherapist (informal)\textquotedblright{}.) Say it once more.
\end{tcolorbox}
\section[(dialogue)]{Retour sur les chapitres 474–481}

\label{lesson:FR-R481-revision-474-481}



\begin{quote}
Say the French for cheerfulness and to laugh.
\end{quote}

\subsection*{Your turn}
Say these aloud:

\begin{itemize}
\raggedright
  \item cheerfulness, to laugh, fit of the giggles, funny, to tickle
  \item jealousy, to envy, grudge, resentful, to reconcile
  \item worry, anxious, to be wary, optimistic, pessimistic
  \item thigh, calf, heel, toe, hip
  \item eyebrow, eyelid, eyelash, nostril, temple
  \item to injure, cut, scratch, burn, bump
  \item fracture, plaster cast, crutch, to limp, bone
  \item sprain, to twist, to swell, splint, physiotherapist
\end{itemize}

\begin{tcolorbox}[breakable,colback=teal!4,colframe=teal!35!black,title={Before you move on}]
Cheerfulness? (\textbf{la gaieté}.) To laugh? (\textbf{rigoler}.) Fit of the giggles? (\textbf{le fou rire}.) Funny? (\textbf{marrant}.) To tickle? (\textbf{chatouiller}.) Jealousy? (\textbf{la jalousie}.) To envy? (\textbf{envier}.) Grudge? (\textbf{la rancune}.) Resentful? (\textbf{rancunier}.) To reconcile? (\textbf{réconcilier}.) Worry? (\textbf{l'inquiétude}.) Anxious? (\textbf{anxieux}.) To be wary? (\textbf{méfier}.) Optimistic? (\textbf{optimiste}.) Pessimistic? (\textbf{pessimiste}.) Thigh? (\textbf{la cuisse}.) Calf? (\textbf{le mollet}.) Heel? (\textbf{le talon}.) Toe? (\textbf{l'orteil}.) Hip? (\textbf{la hanche}.) Eyebrow? (\textbf{le sourcil}.) Eyelid? (\textbf{la paupière}.) Eyelash? (\textbf{le cil}.) Nostril? (\textbf{la narine}.) Temple? (\textbf{la tempe}.) To injure? (\textbf{blesser}.) Cut? (\textbf{la coupure}.) Scratch? (\textbf{l'égratignure}.) Burn? (\textbf{la brûlure}.) Bump? (\textbf{la bosse}.) Fracture? (\textbf{la fracture}.) Plaster cast? (\textbf{le plâtre}.) Crutch? (\textbf{la béquille}.) To limp? (\textbf{boiter}.) Bone? (\textbf{l'os}.) Sprain? (\textbf{l'entorse}.) To twist? (\textbf{tordre}.) To swell? (\textbf{enfler}.) Splint? (\textbf{l'attelle}.) Physiotherapist? (\textbf{le kiné}.)
\end{tcolorbox}
